%%
%% InterruptLab (c) 2021-23 Christopher A. Bohn
%%
%% Licensed under the Apache License, Version 2.0 (the "License");
%% you may not use this file except in compliance with the License.
%% You may obtain a copy of the License at
%%     http://www.apache.org/licenses/LICENSE-2.0
%% Unless required by applicable law or agreed to in writing, software
%% distributed under the License is distributed on an "AS IS" BASIS,
%% WITHOUT WARRANTIES OR CONDITIONS OF ANY KIND, either express or implied.
%% See the License for the specific language governing permissions and
%% limitations under the License.
%%

%%
%% (c) 2021 Christopher A. Bohn
%%

\documentclass[12pt]{article}

\usepackage{fullpage}
\usepackage{fancyhdr}
\usepackage[procnames]{listings}
\usepackage{hyperref}
\usepackage{textcomp}
\usepackage{bold-extra}
\usepackage[dvipsnames]{xcolor}
\usepackage{etoolbox}


% Customize the semester (or quarter) and the course number

\newcommand{\courseterm}{Spring 2022}
\newcommand{\coursenumber}{CSCE 231}

% Customize how a typical lab will be managed;
% you can always use \renewcommand for one-offs

\newcommand{\runtimeenvironment}{your account on the \textit{csce.unl.edu} Linux server}
\newcommand{\filesource}{Canvas or {\footnotesize$\sim$}cse231 on \textit{csce.unl.edu}}
\newcommand{\filesubmission}{Canvas}

% These are placeholder commands and will be renewed in each lab

\newcommand{\labnumber}{}
\newcommand{\labname}{Lab \labnumber\ Assignment}
\newcommand{\shortlabname}{}
\newcommand{\duedate}{}

% Individual or team effort

\newcommand{\individualeffort}{This is an individual-effort project. You may discuss concepts and syntax with other students, but you may discuss solutions only with the professor and the TAs. Sharing code with or copying code from another student or the internet is prohibited.}
\newcommand{\teameffort}{This is a team-effort project. You may discuss concepts and syntax with other students, but you may discuss solutions only with your assigned partner(s), the professor, and the TAs. Sharing code with or copying code from a student who is not on your team, or from the internet, is prohibited.}
\newcommand{\freecollaboration}{In addition to the professor and the TAs, you may freely seek help on this assignment from other students.}
\newcommand{\collaborationrules}{}

% Do you care about software engineering?

\providebool{allowspaghetticode}

\setbool{allowspaghetticode}{false}

\newcommand{\softwareengineeringfrontmatter}{
    \ifboolexpe{not bool{allowspaghetticode}}{
        \section*{No Spaghetti Code Allowed}
        In the interest of keeping your code readable, you may \textit{not} use
        any \lstinline{goto} statements, nor may you use any \lstinline{break}
        statements to exit from a loop, nor may you have any functions
        \lstinline{return} from within a loop.
    }{}
}

\newcommand{\spaghetticodepenalties}[1]{
    \ifboolexpe{not bool{allowspaghetticode}}{
        \penaltyitem{1}{for each \lstinline{goto} statement, \lstinline{break}
            statement used to exit from a loop, or \lstinline{return} statement
            that occurs within a loop.}
    }{}
}

% You shouldn't need to customize these,
% but you can if you like

\lstset{language=C, tabsize=4, upquote=true, basicstyle=\ttfamily}
\newcommand{\function}[1]{\textbf{\lstinline{#1}}}
\setlength{\headsep}{0.7cm}
\hypersetup{colorlinks=true}

\newcommand{\startdocument}{
    \pagestyle{fancy}
    \fancyhf{}
    \lhead{\coursenumber}
    \chead{\ Lab \labnumber: \labname}
    \rhead{\courseterm}
    \cfoot{\shortlabname-\thepage}

	\begin{document}
	\title{\ Lab \labnumber}
	\author{\labname}
	\date{Due: \duedate}
	\maketitle

    \textit{\collaborationrules}
}

\newcommand{\rubricitem}[2]{\item[\underline{\hspace{1cm}} +#1] #2}
\newcommand{\bonusitem}[2]{\item[\underline{\hspace{1cm}} Bonus +#1] #2}
\newcommand{\penaltyitem}[2]{\item[\underline{\hspace{1cm}} -#1] #2}

%%
%% labs/common/semester.tex
%% (c) 2021-22 Christopher A. Bohn
%%
%% Licensed under the Apache License, Version 2.0 (the "License");
%% you may not use this file except in compliance with the License.
%% You may obtain a copy of the License at
%%     http://www.apache.org/licenses/LICENSE-2.0
%% Unless required by applicable law or agreed to in writing, software
%% distributed under the License is distributed on an "AS IS" BASIS,
%% WITHOUT WARRANTIES OR CONDITIONS OF ANY KIND, either express or implied.
%% See the License for the specific language governing permissions and
%% limitations under the License.
%%


% Customize the semester (or quarter) and the course number

\newcommand{\courseterm}{Fall 2022}
\newcommand{\coursenumber}{CSCE 231}

% Customize how a typical lab will be managed;
% you can always use \renewcommand for one-offs

\newcommand{\runtimeenvironment}{your account on the \textit{csce.unl.edu} Linux server}
\newcommand{\filesource}{Canvas or {\footnotesize$\sim$}cse231 on \textit{csce.unl.edu}}
\newcommand{\filesubmission}{Canvas}

% Customize for the I/O lab hardware

\newcommand{\developmentboard}{Arduino Nano}
%\newcommand{\serialprotocol}{SPI}
\newcommand{\serialprotocol}{I2C}
%\newcommand{\displaymodule}{MAX7219digits}
%\newcommand{\displaymodule}{MAX7219matrix}
\newcommand{\displaymodule}{LCD1602}

\setbool{usedisplayfont}{true}

\newcommand{\obtaininghardware}{
    The EE Shop has prepared ``class kits'' for CSCE 231; your class kit costs \$30.
    The EE Shop is located at 122 Scott Engineering Center and is open M-F 7am-4pm. You do not need an appointment.
    You may pay at the window with cash, with a personal check, or with your NCard.
    The EE shop does \textit{not} accept credit cards.
}

% Update to reflect the CS2 course(s) at your institute

\newcommand{\cstwo}{CSCE~156, RAIK~184H, or SOFT~161}

% Do you care about software engineering?

\setbool{allowspaghetticode}{false}

% Which assignments are you using this semester, and when are they due?

\newcommand{\pokerlabnumber}{1}
\newcommand{\pokerlabcollaboration}{
    Sections~\ref{sec:connecting}, \ref{sec:terminology}, \ref{sec:gettingstarted}, \ref{subsec:typesofpokerhands}, and~\ref{subsec:studythecode}: \freecollaboration
    Sections~\ref{sec:completingcard} and~\ref{subsec:completepoker}: \individualeffort
}
\newcommand{\pokerlabdue}{Week of August 29, before the start of your lab section}

\newcommand{\keyboardlabnumber}{2}
\newcommand{\keyboardlabcollaboration}{\individualeffort}
\newcommand{\keyboardlabdue}{Week of January 31, before the start of your lab section}

\newcommand{\pointerlabnumber}{3}
\newcommand{\pointerlabcollaboration}{\individualeffort}
\newcommand{\pointerlabdue}{Week of February 7, before the start of your lab section}

\newcommand{\integerlabnumber}{4}
\newcommand{\integerlabcollaboration}{\individualeffort}
\newcommand{\integerlabdue}{Week of February 14, before the start of your lab section}

\newcommand{\floatlabnumber}{5}
\newcommand{\floatlabcollaboration}{\individualeffort}
\newcommand{\floatlabdue}{soon}

\newcommand{\addressinglabnumber}{6}
\newcommand{\addressinglabcollaboration}{\individualeffort}
\newcommand{\addressinglabdue}{Week of February 28, before the start of your lab section}

%bomblab was 7
%attacklab was 8

\newcommand{\pollinglabnumber}{9}
\newcommand{\pollinglabcollaboration}{\individualeffort}
\newcommand{\pollinglabdue}{Week of April 11, before the start of your lab section}
\newcommand{\pollinglabenvironment}{your \developmentboard-based class hardware kit}

\newcommand{\ioprelabnumber}{\pollinglabnumber-prelab}
\newcommand{\ioprelabcollaboration}{\freecollaboration}
\newcommand{\ioprelabdue}{Before the start of your lab section on April 5 or 6}

\newcommand{\interruptlabnumber}{10}
\newcommand{\interruptlabcollaboration}{\individualeffort}
\newcommand{\interruptlabdue}{Week of April 18, before the start of your lab section}
\newcommand{\interruptlabenvironment}{your \developmentboard-based class hardware kit}

\newcommand{\capstonelab}{ComboLock}    % this will come into play when we generalize capstonelab
\newcommand{\capstonelabnumber}{11}
\newcommand{\capstonelabcollaboration}{\teameffort}
\newcommand{\capstonelabdue}{Week of May 2, Before the start of your lab section\footnote{See Piazza for the due dates of teams with students from different lab sections.}}
\newcommand{\capstonelabenvironment}{your \developmentboard-based class hardware kit}

\newcommand{\memorylabnumber}{12}
\newcommand{\memorylabcollaboration}{This is an individual-effort project. You may discuss the nature of memory technologies and of memory hierarchies with classmates, but you must draw your own conclusions.}
\newcommand{\memorylabdue}{Week of May 2, at the end of your lab section}
\newcommand{\memorylabenvironment}{your \developmentboard-based class hardware kit and your account on the \textit{csce.unl.edu} Linux server}

% Labs not used this semester

\newcommand{\concurrencylabnumber}{XX}
\newcommand{\concurrencylabcollaboration}{\individualeffort}
\newcommand{\concurrencylabdue}{not this semester}

\newcommand{\ssbcwarmupnumber}{XX}
\newcommand{\ssbcwarmupcollaboration}{\freecollaboration}
\newcommand{\ssbcwarmupdue}{not this semester}

\newcommand{\ssbcpollingnumber}{XX}
\newcommand{\ssbcpollingcollaboration}{\individualeffort}
\newcommand{\ssbcpollingdue}{not this semester}

\newcommand{\ssbcinterruptnumber}{XX}
\newcommand{\ssbcinterruptcollaboration}{\individualeffort}
\newcommand{\ssbcinterruptdue}{not this semester}

%%
%% labs/common/storylines.tex
%% (c) 2020-23 Christopher A. Bohn
%%
%% Licensed under the Apache License, Version 2.0 (the "License");
%% you may not use this file except in compliance with the License.
%% You may obtain a copy of the License at
%%     http://www.apache.org/licenses/LICENSE-2.0
%% Unless required by applicable law or agreed to in writing, software
%% distributed under the License is distributed on an "AS IS" BASIS,
%% WITHOUT WARRANTIES OR CONDITIONS OF ANY KIND, either express or implied.
%% See the License for the specific language governing permissions and
%% limitations under the License.
%%

\newcommand{\MeetArchie}{
    \begin{wrapfigure}{r}{0.33\textwidth}
        \centering
        \includegraphics[width=.4\textwidth]{archie}
        \caption{Archie.\\ \footnotesize{Photograph by Bohn.}}
    \end{wrapfigure}

    You're relaxing at your favorite hangout when another customer catches your attention.
    He's rather large (dare I say, \textit{mammoth}), a bit hairy, and looking frustrated in front of his laptop.
    ``I'm Archie,'' he says, ``and I'm trying to teach myself this card game called \textit{Poker}.
    I found this source code that I thought I could use to understand Poker better, but the code is incomplete, and I don't entirely understand what's there.
    Could you explain the code to me, please?'
}

\newcommand{\GetHired}{
    Archie's face lights up in a very big smile.
    ``Thanks!''
    After pausing in thought for a moment, he says, ``Say, I've got a new startup company that could really use your help.
    Are you interested?
    It'll be exciting!''
}

\newcommand{\FirstDayOnTheJob}{
    \begin{wrapfigure}{r}{0.33\textwidth}
        \centering
        \includegraphics[width=.4\textwidth]{some-expenses-spared}
        \caption{Some expenses were spared.\\ \footnotesize{Original images \textcopyright\ Universal Studios and Amblin Entertainment, Inc. Meme creator unknown.}}
    \end{wrapfigure}

    You've recently been hired to help get the Pleistocene Petting Zoo get started.
    Your new employer, Archie, is surprisingly honest: he admits to you that some expenses were spared.
    Archie cheerfully points out that any challenge is also an opportunity to succeed.
    You suspect your job will offer plenty of ``opportunities to succeed.''
}

\newcommand{\HasKeyboard}{
    Great news!
    Archie brings you your new keyboard.
    He also brings you a problem of his own.
    Because you were held up with the broken keyboard, Archie decided to try some programming on his own, and his code is behaving strangely.
}

\newcommand{\ArchieWroteSmellyCode}{
    Working at the Pleistocene Petting Zoo certainly is proving to be interesting.
    You're glad that you don't have to worry about the problem of the giant sloths very slowly chasing their handlers, but now it seems that Archie has decided to try to write a program or two.
    At a glance, his code is smellier than the wooly rhinoceros' enclosure.
    But you take a closer look anyway to try to understand why his code acts strangely.
}

\newcommand{\InsurancePreview}{
    You hear somebody enter the room.
    ``\textit{Frankenstein}, `boat','' is the challenge, and she answers, ``borne.''
    Archie introduces you to the new arrival, ``Lil, this is our new developer, the one who wrote the app we just used.''
    He turns to you: ``This is Lilith Redd from business operations.''
    He turns back to her and continues, ``Lil, what's the good word?''

    ``The word isn't good, I'm afraid.
    I just heard back from the insurance company.''
}

\newcommand{\OnLoanToEclecticElectronics}{
    All work at the Pleistocene Petting Zoo has stopped while Archie tries to find a $\cancelto{\mathrm{reasonable}}{\mathrm{gullible}}$ insurance company.
    Rather than furloughing staff, he's asked everybody to help out with his other startup companies for a week or two.
    He specifically asked that you help out with Eclectic Electronics.

    Herb Bee, the chief engineer, explains that Eclectic Electronics is developing a patent-pending C-licon tool that will convert C code into an integrated circuit that has the same functionality as the original C code.
    To test it out, he tasked you with writing the code to implement an Arithmetic Logic Unit (ALU).
    Your task will be to implement integer addition, subtraction, multiplication, and division.
    Even though high-level languages' \textit{logical} boolean operations normally are not part of an ALU, Herb wants you to include these in the ALU to see if that can make some programs run faster.
    Because bitwise operations and bit shift operations have been implemented, you will be able to use C's bitwise and bit shift operators, but because arithmetic operations have not yet been implemented, you cannot use C's arithmetic operators.
    Because C library functions generally make use of arithmetic operations (which have not yet been implemented), you cannot use library functions.
}

\newcommand{\SuccessfulALU}{
    Herb smiles as he hands you the the test results from the latest integrated circuit fab batch.
    ``C-licon successfully turned your code into an ALU.
    Nicely done!''
    I think maybe it's time to use C-licon to see if we can improve the Floating Point Unit (FPU) on our experimental microprocessor.
}

\newcommand{\WriteAnFPU}{
    Herb tells you that, Eclectic Electronics tested the integrated circuit that the C-licon tool created from your ALU code, and they've concluded that C-licon is ready to use for their new experimental microprocessor.
    He tasks you with writing C code (that will be used by the C-licon tool) to implement a Floating Point Unit (FPU).
    Your task will be to implement floating point addition, subtraction, multiplication, and division.
    You can use any bit operations and, thanks to the ALU you wrote, you can use any integer arithmetic operations (use the conventional + - * / operators).
    Because the FPU has not yet been implemented, you cannot use C's floating point operations, you cannot use \lstinline{float}s nor \lstinline{double}s, and you cannot use library functions.
}

\newcommand{\GoingBackToTheZoo}{
    Lil enters the room.
    Herb challenges her: ``\textit{Gulliver's Travels}, `endian','' and Lil answers, ``ends.''

    Lil walks up to you and says, ``We have the insurance situation taken care of, and it's time to get the Zoo ready for guests.
    We're reassembling the tech team, and there's plenty of work to do.''

    You smile.
    ``That's good news!''

    Lil's face is hard to read.
    ``Well, yes and no.
    It's good that you'll be able to resume work on the Zoo's systems.
    But while Archie was waiting for us to fix the insurance situation, he got bored and -- cutting a long story short -- he ended up creating some new `opportunities' that we need you `to succeed' at.''
}

\newcommand{\SettledIntoRoutine}{
    You've settled into a comfortable routine at the Pleistocene Petting Zoo.
    While your job isn't quite as exciting as that of the saber-toothed tigers' dentist, it still has something new and interesting almost every day.

    Archie announces that he heard that hand-crafted assembly code can be faster than high-level language code.
    You try to explain that while this may have been true decades ago, modern optimizing compilers generate code faster than what a typical programmer can achieve with assembly code.
    Archie doesn't believe you and insists that you write the zoo's new cipher program in x86 assembly code.
}

\newcommand{\NewmanRanOffWithSamples}{
    Archie is hurriedly packing is trunk, like he's about to leave on a short-notice urgent trip.
    Before charging out the door, he pauses to tell you, ``Newman just stole some of our samples.
    I need to track him down before he sells them to the Supersized Safari Syndicate.
    I guess this means you're in charge of the Zoo's computer system now.
    Don't worry, you'll be fine. What could possibly go wrong?''
}

\newcommand{\BombLabIntroduction}{  % Ties Bryant & O'Halloran's Bomb Lab into the Pleistocene Petting Zoo story
    In a jarring collision of movie franchises, the CEO of Virtucon makes a Zoom call to the Pleistocene Petting Zoo.
    For some reason that nobody really explains, you're the only person available to handle the situation.
    The guy, who sounds kind of like an animated ogre, demands that the Pleistocene Petting Zoo deliver to him a megalodon shark with a head-mounted laser capable of emitting a beam of pure antimatter.

    You blurt out, ``Then it's not a laser,'' and then try to explain to him that megalodons are from the Miocene epoch, and expecting to find them at the Pleistocene Petting Zoo would be as ridiculous as a Cretaceous-period tyrannosaur at a Jurassic-themed park.

    ``Zip it!'' commands the guy who kind of looks like the host of a public-access show you used to watch.
    ``Since you won't meet my demand, my minions have placed a `binary bomb' under your zoo.
    Because I like really convoluted plans, we put software on your Linux server that controls the bomb.
    If you do nothing, the bomb will explode.
    If you turn off the Linux server, the bomb will explode.
    If you go slower than 50mph, the bomb will -- no, never mind that last part.

    ``The bomb software consists of a sequence of phases.
    Each phase expects you to type a particular string on \texttt{stdin}.
    If you type the correct string, then the phase is {\em defused} and the bomb proceeds to the next phase. Otherwise, the bomb {\em explodes}.
    The bomb is defused when every phase has been defused.

    ``Your mission, which you have no choice but to accept, is to defuse your bomb before the due date.
    Good luck, and welcome to the bomb squad!''
}

\newcommand{\FoodLockersAreStuck}{
    Having saved the Zoo from Dr. Evil's binary bomb, you relax back in your chair and think about taking a break.
%SPRINGBREAK
    Maybe an entire week in which you don't have solve any problems or meet any deadlines -- that'd be real nice.
%FALLBREAK
    % Maybe 4-day weekend in which you don't have solve any problems or meet any deadlines -- that'd be real nice.

    Another Zoom call comes in.
    \textit{What now!?} you wonder as you take your feet off of the desk to answer the call.
    An uncomfortable-looking animal handler says, ``We can't unlock the food lockers.
    It's the animals' feeding time, and we can't open the food lockers!
    It's feeding time, we can't get to the animals' food, and,'' his eyes dart nervously toward the animal enclosures, ``and many of them have sharp, pointy claws and others have big, stompy feet.''
}

\newcommand{\AttackLabIntroduction}{    % Ties Bryant & O'Halloran's Attack Lab into the Pleistocene Petting Zoo story
    You managed to keep the Pleistocene Petting Zoo from blowing to smithereens, but it turns out that Dr. Evil's minions weren't too careful when they put the bomb control software on the Zoo's Linux server.
    The software that controls the food locker has been heavily damaged!
    The functions that unlock the food locker doors are still present, but there's no way to activate those functions.

    You then recall what Archie told you when he hired you: some expenses were spared.
    You run the machine code through a disassembler and quickly see that it has a buffer overflow vulnerability.
    Before the situation in the dire wolf enclosure gets too dire, you sit down and get to work.

    The \function{ctarget} code runs on an older machine that allows executable code to be present on the stack, so it's vulnerable to a conventional code injection buffer overflow attack.
    \begin{itemize}
        \item Phase 1 (\function{touch1}) unlocks the food locker so the animal handlers can prepare the food.
        \item Phase 2 (\function{touch2}) opens the doors between the food locker and the carnivore enclosures;
        you will need to pass a cookie to the function to authenticate yourself.
        \item Phase 3 (\function{touch3}) closes the doors between the food locker and the carnivore enclosures.
    \end{itemize}
    The \function{rtarget} code runs on a newer machine that does not allow executable code to be present on the stack, so you'll have to conduct a return-oriented programming attack on it.
    \begin{itemize}
        \item Phase 4 (\function{touch2)} opens the doors between the food locker and the herbivore enclosures;
        you will need to pass a cookie to the function to authenticate yourself.
        \item Phase 5 (\function{touch3}) closes the doors between the food locker and the herbivore enclosures.
    \end{itemize}
}

\newcommand{\MostAnimalsAreFed}{
    Before you take on the Phase 5, pause to consider what you have accomplished so far.
    In Phases 2 and 3, you caused a program to execute machine code of your own design.
    If {\sc ctarget} had been a network server, you could have injected your own code into a distant machine.
    In Phase 4, you circumvented two of the main devices modern systems use to thwart buffer overflow attacks.
    Although you did not inject your own code, you were able inject a type of program that operates by stitching together sequences of existing code.
    Also, all animals have been fed, the carnivores are still in their enclosure, the mammoths can't fit through the herbivore door, and only the giant sloths seem interested in very slowly escaping.
}

\newcommand{\ArchieReturns}{
    Archie returns from tracking down Newman, who'd run off with some of the Pleistocene Petting Zoo's samples shortly before Dr.~Evil's Zoom call.
    ``It turns out he didn't get very far at all,'' Archie sighs.
    ``He ran into a flock of terror birds as he was leaving, and we found him in one of the emergency shelters.

    Archie smiles. ``I trust things were uneventful while I was away?''
}

\newcommand{\PickingUpNewmansProject}{
    Archie seems genuinely surprised that Newman is refusing to go back to work.
    ``You would think that he'd be grateful for being rescued from that flock of terror birds.''
    Before you can wonder out-loud whether it would be a good idea to trust someone who had just tried to sell trade secrets to a competitor, Archie gives you your new task.

    ``Because Newman isn't cooperating, I need you to finish the project he was working on.
    As you can imagine, duplicating the genetic information for our exhibits can take a long time, and Newman realized that we might be able to duplicate the data faster if we had a concurrent program which has one thread reading from the original data and another thread writing the copy.
    Unfortunately, he ran off to sell samples to the  Supersized Safari Syndicate before finishing the duplicator.
    Right now the duplicator seems to work, but it usually makes imperfect copies.
    Have you ever seen a paleolama with two noses, four eyes, and no ears!?''
}

\newcommand{\WeNeedBetterDetection}{
    Between Newman trying to sell samples to a competitor, that weird guy almost blowing up the zoo, and the animals almost escaping, Archie is getting worried.
    ``I think we need to introduce additional protective measures.
    As useful as your challenge-response app is in helping us detect intruders, I think it's now clear that we need something that will detect someone -- or some\textit{thing} -- when they're someplace they shouldn't be, even when no one else is around.
    I've asked the team at Eclectic Electronics to put something together.''
}

\newcommand{\WeNeedBetterLocks}{
    Between Newman trying to sell samples to a competitor, that weird guy almost blowing up the zoo, and the animals almost escaping, Archie is getting worried.
    ``I think we need to introduce additional protective measures.
    As useful as your challenge-response app is in helping us detect intruders, I think it's now clear that we need something that will keep someone -- or some\textit{thing} -- out of places they shouldn't be, even when no one else is around.
    I've asked the team at Eclectic Electronics to put something together.''
}

\newcommand{\IntroduceHardware}{
    Archie walks up to you, along with Herb Bee from Eclectic Electronics.
    Herb is holding a tangled mess of electronics.
    Archie explains, ``Herb here has developed a prototype of a device that he thinks will be useful for our physical security needs, as well as a few other applications around here. He calls it the \textit{Cow Pi}.''

% TODO: parameterize based on which microcontroller is actually being used
    You look at the device in Herb's hands and see the \nano\ central to the circuit.
    ``Isn't \textit{-Pi} typically used as a suffix for circuits that use a Raspberry Pi instead of an Arduino?''

    Herb replies, ``Typically, yes, but \textit{Cowduino} isn't very punny, is it?''

    Archie chimes in, ``Maybe with the right emphasis: \textit{Cow-DOO-ino}.''

    ``That's kind of subtle, don't you think? How will people know to put the emPHAsis on that sylLAble?''

    ``I think we're getting off topic here,'' you point out.
    ``How can I help?''

    ``Oh, right,'' Herb says, ``We'd like you to kick its proverbial tires.
    Let's start off with something simple, like a number builder tool.''
}

\newcommand{\JeffGoldblum}{
    Herb looks over your work.
    ``Hmm, yes. I think this is coming along nicely.
    Let's run a few more tests.''

    Archie storms into the room.
    ``We have \textit{got} to do something about security!
    How's that doodad coming along?
    Because there's now a half-man/half-fly in the labs going on-and-on about Chaos Theory and how if we just give him a MacBook and a spaceship then he'll be able to get the Lord of Thunder to travel across the 8th Dimension.
    Is that thing just about ready?''

    Herb shakes his head, ``No, not quite yet. It should be ready in about a week.''
}

\newcommand{\DisdainfulHerb}{
    Smoke wafts from Herb's soldering iron as he looks up when you approach.
    Cleaning the iron's tip, he quotes:
    ``Somebody once said, `The three most dangerous things in the world are a programmer with a soldering iron, a hardware engineer with a software patch, and'{}'' -- he glances nervously in Archie's direction -- ``{}`a user with an idea.'\footnote{
        Rick Cook, \textit{The Wizardry Consulted}, 1995.
    }$^{\mathrm{,}}$\footnote{
        The notion of being wary of programmers wielding screwdrivers or soldering irons long pre-dated this quote, as there are apocryphal tales of people who found it easier to modify the hardware to suit the software rather than the other way around.
    }''
}

\newcommand{\NumberConversionTool}{     % Since we're now allowing `sprintf()` with the LCD1602, converting between decimal and hexadecimal is trivial; it still might be okay for 7-segment displays
    Herb gets straight to the point.
    ``We promised Archie that we'd be able to start using the Cow Pi to build systems in a week.
    So far we've tested its input/output functionality, but we still need to test its timer and also whether we can take inputs without constantly polling the input devices.
    As before, we don't need to do anything too fancy;
    let's try a number base conversion tool.''
}

\newcommand{\LessDisdainfulHerb}{
    Smoke wafts from Herb's soldering iron as he looks up when you approach.
    Cleaning the iron's tip, he notes:
    ``Somebody once said that one of the most dangerous things in the world is a programmer with a soldering iron.''\footnote{
        ``The three most dangerous things in the world are a programmer with a soldering iron, a hardware engineer with a software patch, and a user with an idea.'' -- Rick Cook, \textit{The Wizardry Consulted}, 1995.
    }$^{\mathrm{,}}$\footnote{
        The notion of being wary of programmers wielding screwdrivers or soldering irons long pre-dated this quote, as there are apocryphal tales of people who found it easier to modify the hardware to suit the software rather than the other way around.
    }
}

\newcommand{\RemoteControlledCar}{

    About this time, Archie walks by, thinking about electric carts to transport visitors around the Pleistocene Petting Zoo.
    ``They probably should be remote-controlled.''
    He looks at you and Herb, and asks, ``Do you think you could make a cart a remote-controlled cart?''

    You ask the obvious question, ``Are there carts here already?''

    Archie waves his hand in the air, dismissing that detail, ``Not yet, but could you make the remote-control?''
    
    You hestatingly summarize: ``You want a cartless remote-controlled cart?''

    Archie beamingly smiles, ``Exactly!''

    Herb jumps in, ``Yes, we'll do it.''
    Herb looks at you and adds, ``It'll give us a chance to test the Cow Pi's timer and whether we can take inputs without constantly polling the input devices.''
}

\newcommand{\LauraDern}{
    You and Herb look for Archie in the Pleistocene Petting Zoo's labs to give him the good news, and you find a blond woman wearing cargo shorts, butchering a Gilbert and Sullivan song\dots \\ \\
    \textmusicalnote\ I am the very model of a modern vice admiral \textmusicalnote \\
    \textmusicalnote\ I've information about all things paleobotanical \textmusicalnote \\
    \textmusicalnote\ And I've been up to my armpits in problems scatological \textmusicalnote \\
    \textmusicalnote\ During the regency I had experience matriarchical \textmusicalnote \\
    \textmusicalnote\ I plot space travel, normal and superluminal \textmusicalnote \\
    \textmusicalnote\ (Even if I challenge the Pauli exclusion principle) \textmusicalnote \\

    ``I don't know how these people keep getting into our labs.
    \textit{Please} tell me that you have good news,'' pleads Archie.

    ``Yes, the Cow Pi is ready for whatever you need: calculators, security systems, parking meters -- you name it,'' Herb cheerfully responds.

    ``Excellent.''
    Archie turns to you.
    ``I'd like you and Newm... no, \textit{not} Newman.
    I'd like you and someone else on the staff to get started right away.
    Here's what I'd like to have built first.''
}

\newcommand{\CalculatorNeeded}{
    ``I have various teams working on different projects around here to improve security,'' Archie reminds you.
    He glances toward the Zoo's labs, where there's now a guy who looks like the actor who portrayed the fictional actor who portrayed the Norse god Odin, trying to avoid children while wistfully talking about raising rabbits in Montana.
    You briefly wonder why there are children someplace where there are also carnivorous megafauna, and then you remember that you work at a petting zoo.
    ``What I need your team to do,'' Archie continues, ``is make a four-function calculator so that we can quickly and easily determine whether we have the correct number of specimens, or if any are missing.''
}

\newcommand{\CalculatorCounting}{
    Technicians are using your calculator to compute how many specimens are still present in the lab, and establish that all specimens are accounted for after Newman's attempted theft.
    As reports come in of facilities getting secured with Cow Pi-based locks and passages being monitored with Cow Pi-based motion sensors, Archie smiles and tells you that this was a job well done.
    With all of the excitement neatly wrapped-up and arriving at a satisfactory conclusion, you look forward to a boring career in which there's absolutely no screaming and running for your life.
}

\newcommand{\CombinationLockNeeded}{
    ``I have various teams working on different projects around here to improve security,'' Archie reminds you.
    He glances toward the Zoo's labs, where there's now a guy who looks like the actor who portrayed the fictional actor who portrayed the Norse god Odin, trying to avoid children while wistfully talking about raising rabbits in Montana.
    You briefly wonder why there are children someplace where there are also carnivorous megafauna, and then you remember that you work at a petting zoo.
    ``What I need your team to do,'' Archie continues, ``is make a combination lock so that only authorized people can get into our lab facilities.''
}

\newcommand{\CombinationLockInstalled}{
    After fastening the new electronic combination lock to the lab door, Archie smiles and tells you that this was a job well done.
    With all of the excitement neatly wrapped-up and arriving at a satisfactory conclusion, you look forward to a boring career in which there's absolutely no screaming and running for your life.
}

\newcommand{\RangeFinderNeeded}{
    ``I have various teams working on different projects around here to improve security,'' Archie reminds you.
    He glances toward the Zoo's labs, where there's now a guy who looks like the actor who portrayed the fictional actor who portrayed the Norse god Odin, trying to avoid children while wistfully talking about raising rabbits in Montana.
    You briefly wonder why there are children someplace where there are also carnivorous megafauna, and then you remember that you work at a petting zoo.
    ``What I need your team to do,'' Archie continues, ``is make a range finder that will alert us when someone -- or some\textit{thing} -- gets too close to someplace they shouldn't be.''
}

\newcommand{\RangeFinderDetecting} {
    A technician installing a new range finder outside the lab door briefly sets off the alarm, but then the range finder falls quiet and faithfully reports that nothing is approaching.
    As reports come in of facilities getting secured with Cow Pi-based locks, and of accurate speciment counts accomplished with Cow Pi-based calculators, Archie smiles and tells you that this was a job well done.
    With all of the excitement neatly wrapped-up and arriving at a satisfactory conclusion, you look forward to a boring career in which there's absolutely no screaming and running for your life.
}

\usepackage{amsmath}
%\usepackage{array,color,colortbl}
\definecolor{LightGreen}{rgb}{0.88,1,0.88}
\usepackage{multicol}
\usepackage{CJKutf8}
\usepackage{gensymb}
\setlength{\columnsep}{2cm}


% \captionsetup{width=.8\linewidth}

% \lstset{language=c, numbers=left, showstringspaces=false,
%     moredelim = [s][\ttfamily]{/*}{*/} % I shouldn't need this parameter!
%     }

\renewcommand{\labnumber}{\interruptlabnumber}
\renewcommand{\labname}{Using Interrupt-Driven Input/Output}
\renewcommand{\shortlabname}{interrupt-driven i/o -- interruptlab}
\renewcommand{\collaborationrules}{\interruptlabcollaboration}
\renewcommand{\duedate}{\interruptlabdue}
\newcommand{\nano}{\developmentboard} % TODO: replace \nano with \developmentboard
\renewcommand{\runtimeenvironment}{\interruptlabenvironment}
\pagelayout
\begin{document}
    \labidentifier

    \pdfbookmark[1]{Frontmatter}{frontmatter}                   In this assignment, you will write code for \runtimeenvironment\ that will use new electronic devices to interact with the physical world.

The instructions are written assuming you will edit the code in the Arduino IDE and run it on \runtimeenvironment, constructed according to the pre-lab instructions.
If you wish, you may edit the code in a different environment; however, our ability to provide support for problems with other IDEs is limited.

\section*{Learning Objectives}

After successful completion of this assignment, students will be able to:
\begin{itemize}
    \item Work collaboratively on a hardware/software project
    \item Design and implement a simple embedded system
    \item Expand their programming knowledge by consulting documentation
\end{itemize}

\subsection*{Continuing Forward}

This penultimate lab assignment does not contribute to the final lab assignment.
By integrating elements of what you learned in this course, and by demonstrating that you can review documentation to learn on your own, to design a small embedded system, you will show how much progress you have made this semester.

\section*{During Lab Time}

During your lab period, coordinate with your group partner(s) to decide on your working arrangements.
Unless you're only going to work on the assignment when you're together, you may want to set up a private Git repository that is shared with your partner(s).
With your partner(s), modify your hardware kit as described in Section~\ref{sec:hardwareMods}.
Then, think through your system's design and begin implementing it.
The TAs will be available for questions.


    \softwareengineeringfrontmatter

    \section*{Scenario}                                         %\DisdainfulHerb \NumberConversionTool
                                                                \LessDisdainfulHerb \RemoteControlledCar

    \section{Assignment Summary}                                This assignment is principally about getting comfortable when explicitly working with memory.
Being able to think about a value and a reference to that value distinctly will improve your programming skills in any language.

Before you do so, in Section~\ref{sec:archiesCode} you will examine Archie's code.
Parts of Archie's programs use code that the C standard explicitly states will result in undefined behavior.
By understanding the mistakes that Archie made, we hope that you can avoid them in your own code.

In Section~\ref{sec:challengeResponse}, you will build and use a linked list.
This will require you to allocate space for the list's nodes and manipulate pointers that connect the nodes to each other.

\ifboolexpe{not bool{allowspaghetticode}}{
    There are no particular restrictions in this assignment other than those common to most lab assignments in this course.
    You can check whether you're using a \lstinline{goto} or \lstinline{continue} statement, or whether you're using \lstinline{break} or \lstinline{return} to exit a loop, by running the constraint-checking Python script:
    \texttt{python constraint-check.py linkedlistlab.json}
}{}


    \section{Getting Started} \label{sec:GettingStarted}        Download \textit{\shortlabname.zip} or \textit{\shortlabname.tar} from \filesource\ and copy it to \runtimeenvironment.
Once copied, unpackage the file.
Four of the five files (\textit{alu.h}, \textit{basetwo.c}, \textit{alu.c}, and \textit{integerlab.c}) contain the starter code for this assignment.
The last file (\textit{Makefile}) tells the \texttt{make} utility how to compile the code.
To compile the program, type:

\texttt{make}

This will produce an executable file called \textit{integerlab}.

When you run the program with the command \texttt{\textbf{\textit{./integerlab}}}, you will be prompted:

\begin{verbatim}
    Enter a one- or two-operand logical expression,
        a two-operand comparison expression, a two-operand arithmetic expression,
        "lg <value>" or "exponentiate <value>" to test your powers-of-two code,
        "is_negative <value>" to determine if 2's complement value is negative,
        "add1 <binary_value1> <binary_value2> <carry_in>" for 1-bit full adder,
        "add32 <hex_value1> <hex_value2> <carry_in>" for 32-bit ripple-carry adder,
        or "quit":
\end{verbatim}

When you enter a value, if it is prepended with \texttt{\textbf{\textit{0x}}} then the parser will parse it as a hexadecimal value;
otherwise, except as noted in Sections~\ref{subsec:one-bit-full-adder} and \ref{subsec:ripple-carry-adder}, the parser will treat it as a decimal value.

For now, type \texttt{\textbf{\textit{quit}}} to exit the program.

\subsection{Description of IntegerLab Files}

\subsubsection{integerlab.c}

Do not edit \textit{integerlab.c}.

This file contains the driver code for the lab.
It parses your input, calls the appropriate arithmetic function, and displays the output.

\subsubsection{alu.h} \label{subsubsec:alu.h}

Do not edit \textit{alu.h}.

This header file contains two type definitions:
\begin{description}
    \item[one\_bit\_adder\_t] is a structure to hold the 1-bin inputs (\lstinline{a}, \lstinline{b}, \lstinline{c_in}) and 1-bit outputs (\lstinline{sum}, \lstinline{c_out}) of a one-bit full adder.
    \item[alu\_result\_t] is a structure to hold the outputs from an arithmetic logic unit.
        Its fields are:
        \begin{itemize}
            \item \lstinline{result}, a 16-bit bit vector that is considered ``the'' result of the computation
            \item \lstinline{supplemental_result}, a 16-bit bit vector that stores additional result data from instructions that place their results in two registers
            \item \lstinline{unsigned_overflow}, a 1-bit flag to indicate whether overflow occurred when interpreting the source operands as unsigned values
            \item \lstinline{signed_overflow}, a 1-bit flag to indicate whether overflow occurred when interpreting the source operands as signed values
            \item \lstinline{divide_by_zero}, a 1-bit flag to indicate whether there was an attempt to divide by zero.
        \end{itemize}
\end{description}

The header file also contains two macros, \function{is_zero()} and \function{is_not_zero()} to bootstrap your ALU code.
These macros act like functions and return a boolean value to indicate whether an integer is 0 or not.\footnote{
    The astute student will quickly realize that \function{is_not_zero()} is not necessary and, with a little thought, will realize that they can \function{is_zero()} as a function within the constraints of this assignment.}

The header file also contains several function declarations.
The requirements for these functions will be discussed later in this assignment.

\subsubsection{basetwo.c}

This is the first of two files that you will edit.

There are two functions in \textit{basetwo.c} that will allow you to demonstrate an understanding of powers-of-two and/or an understanding of some uses of bit shifts.
\begin{description}
    \item[lg()] returns the base-2 logarithm of its argument, assuming its argument is a positive power-of-two;
        if the argument is 0 or is not a power-of-two, then there are no guarantees about the function's return value
    \item[exponentiate()] creates a power-of-two by raising 2 to the provided exponent, assuming the exponent is a non-negative value strictly less than 32;
        if the argument is negative or is greater than 31, then there are no guarantees about the function's return value
\end{description}
These functions are inverses of each other: $x = \log_2 2^x$, and $y = 2^{\log_2 y}$.

Strictly speaking, you can write your ALU code without these functions;
however, some students in the past had difficulty finding solutions for their ALU code without obtaining a base-2 logarithm and/or calling a function to create a power-of-two.
Rather than tempt you to violate one of the assignment's constraints by calling the \textit{math} library's \function{log2()}, \function{exp2()}, and/or \function{pow()} functions, we now have you write your own code for these functions.

\subsubsection{alu.c}

This file will contain most of the code that you write, and the functions in \textit{alu.c} are in the order in which you will likely write them.
\begin{itemize}
    \item A simple check
        \begin{description}
            \item[is\_negative()] returns a boolean value to indicate whether the argument, when interpreted as a signed integer, is negative
        \end{description}
    \item Equality comparisons
        \begin{description}
            \item[equal()] returns \lstinline{true} if and only if $value1 = value2$
            \item[not\_equal()] returns \lstinline{true} if and only if $value1 \not = value2$
        \end{description}
    \item Logical operations
        \begin{description}
            \item[logical\_not()] returns the logical inverse of the argument
            \item[logical\_and()] returns the logical conjunction of the two arguments
            \item[logical\_or()] returns the logical disjunction of the two arguments
        \end{description}
    \item Addition and subtraction
        \begin{description}
            \item[one\_bit\_full\_addition()] performs addition for one bit position, determining both the sum bit and the carry-out bit
            \item[ripple\_carry\_addition()] adds two 32-bit values to each other and to a carry-in bit
            \item[add()] adds two 16-bit values to each other
            \item[subtract()] subtracts a 16-bit value from another
        \end{description}
    \item Inequality comparisons
        \begin{description}
            \item[less\_than()] returns \lstinline{true} if and only if $value1 < value2$
            \item[at\_most()] returns \lstinline{true} if and only if $value1 \leq value2$
            \item[at\_least()] returns \lstinline{true} if and only if $value1 \geq value2$
            \item[greater\_than()] returns \lstinline{true} if and only if $value1 > value2$
        \end{description}
    \item Unsigned multiplication and division
        \begin{description}
            \item[multiply\_by\_power\_of\_two()] multiplies the first argument by the second, assuming that the second argument is zero or a power of two;
                there are no guarantees if this assumption is not satisfied
            \item[unsigned\_multiply()] multiplies two 16-bit values to each other, if the arguments are interpreted as unsigned integers
            \item[unsigned\_divide()] divides a 16-bit value by another, if the arguments are interpreted as unsigned integers
        \end{description}
    \item Signed multiplication and division (bonus credit)
    \begin{description}
        \item[signed\_multiply()] multiplies two 16-bit values to each other, if the arguments are interpreted as signed integers
        \item[signed\_divide()] divides a 16-bit value by another, if the arguments are interpreted as signed integers
    \end{description}
\end{itemize}


    \section{Cart Controller Specification} \label{sec:spec}    %! suppress = LabelConvention
\begin{enumerate}
    \item Definitions
        \begin{description}
            \item[Threshold Range] The distance at which a detected object will cause the system to produce an audible alarm
            \item[Strobe] An illumination of an LED for 50ms
            \item[Chirp] A 5kHz tone lasting 50ms
        \end{description}
    \item \label{spec:modes} The slide-switches shall control the mode of operation.
        \begin{enumerate}
            \item Both switches in the \textit{left position}: Normal Operation
            \item The \textbf{left switch} in the \textit{right position} and the \textbf{right switch} in the \textit{left position}: Single-Pulse Operation
            \item Both switches in the \textit{right position}: Threshold Adjustment
            \item The \textbf{left switch} in the \textit{left position} and the \textbf{right switch} in the \textit{right position}: Continuous Tone
        \end{enumerate}
    \item \label{spec:continuousTone} When the range finder and alarm system is in \textbf{Continuous Tone} mode:
        \begin{enumerate}
            \item The system shall not detect the range of nearby objects: it shall neither emit nor detect ultrasound
            \item The system shall not illuminate an LED
            \item The system shall produce a continuous 5kHz audible tone
        \end{enumerate}
    \item \label{spec:thresholdAdjustment} When the range finder and alarm system is in \textbf{Threshold Adjustment} mode:
        \begin{enumerate}
            \item The system shall not detect the range of nearby objects: it shall neither emit nor detect ultrasound
            \item The system shall not produce an audible tone
            \item The system shall not illuminate an LED
            \item The system shall prompt the user on the \textbf{display module} to enter the threshold range, in centimeters
                \begin{itemize}
                    \item You may assume that the user understands what the system does: the prompt must be \textit{meaningful}, but it may be \textit{succinct}
                \end{itemize}
            \item The user shall be able to enter the threshold range, in centimeters, using the numeric keypad
                \begin{enumerate}
                    \item The user will enter the range in decimal
                    \item The system shall echo the user's input, digit-by-digit, on the \textbf{display module} as they type their input
                    \item The user will indicate that they have finished entering their input by pressing the `\#' key
                \end{enumerate}
            \item After the user has completed their input, then any value less than 50cm or greater than 400cm shall be rejected as an invalid threshold;
                the system shall display a helpful error message on the \textbf{display module} and re-prompt the user to enter the threshold range
            \item After the user has completed their input, and if the input is valid, then the system shall display ``Threshold (value)cm'' on the \textbf{display module}, where ``(value)'' shall be the numeric threshold range in centimeters;
                this shall be displayed until the user takes the system out of the Threshold Adjustment mode
        \end{enumerate}
    \item \label{spec:singlePulseOperation} When the range finder and alarm system is in \textbf{Single-Pulse Operation} mode: \\
        {\footnotesize Note: the references here to ``the pushbutton'' refer to the Cow Pi's original \textbf{left pushbutton}}
        \begin{enumerate}
            \item The system shall neither emit nor receive ultrasound, shall not produce an audible tone, and shall not illuminate an LED, \textit{until} the pushbutton has been pressed
            \item Whenever the user presses the \textbf{pushbutton}, the system shall emit one ultrasonic pulse
            \item If the system does not receive that pulse's echo, the system shall resume waiting for the user to press the pushbutton
            \item \label{spec:singlePulseResponse} If the system does receive that pulse's echo:
                \begin{enumerate}
                    \item The system shall strobe both LEDs once
                    \item If the object's distance is less than the threshold range then the system shall emit one chirp
                    \item The system shall then resume waiting for the user to press the pushbutton
                \end{enumerate}
        \end{enumerate}
    \item \label{spec:normalOperation} When the range finder and alarm system is in \textbf{Normal Operation} mode, the system shall repeatedly:
        \begin{enumerate}
            \item \label{spec:emitUltrasound} Emit an ultrasonic pulse and take no further action until either the system receives an echo or the system can establish that it will not receive an echo
            \item If the system can establish that it will not receive an echo:
                \begin{enumerate}
                    \item The system shall stop strobing the LEDs and chirping if it previously was doing so
                    \item The system shall repeat the action in Requirement~\ref{spec:emitUltrasound}
                        {\footnotesize
                        \begin{itemize}
                             \item Repeating the action may be postponed until after necessary quiescent periods
                        \end{itemize}}
                \end{enumerate}
            \item If the system receives an echo:
                \begin{enumerate}
                    \item The system shall compute and display the object's distance, in centimeters
                    \item The system shall compute and display the object's rate of approach, in centimeters per second
                    \item The system shall strobe both LEDs at the rate described by Table~\ref{tab:alarmPeriods}
                    \item If the object's distance is less than the threshold range, the system shall emit chirps at the rate described by Table~\ref{tab:alarmPeriods}
                    \item After the system has completed the speed and distance calculations, it shall repeat the action in Requirement~\ref{spec:emitUltrasound}
                        {\footnotesize
                        \begin{itemize}
                            \item Repeating the action may be postponed until after necessary quiescent periods
                        \end{itemize}}
                \end{enumerate}
        \end{enumerate}
    \item All mechanical inputs shall be properly debounced
    \item The system shall always be responsive to user input
        \begin{itemize}
            \item The user will never press two keys on the numeric keypad at the same time
            \item The system may ignore changes to the slide-switches while the user enters a new threshold range
            \item The system may block while displaying an error message
            \item But for those exceptions, there shall be no noticeable lag when responding to an input
        \end{itemize}
\end{enumerate}

\begin{table}
    \centering
    \begin{tabular}{||c|c|c|c||} \hline\hline
        \multirow{2}{*}{distance}   & LEDs on /         & LED off /         & total     \\
                                    & Piezo sounding    & Piezo silent      & period    \\ \hline\hline
        $distance \geq 250cm$          & $50ms$            & $2450ms$          & $2500ms$  \\ \hline
        $200cm \leq distance < 250cm$  & $50ms$            & $1950ms$          & $2000ms$  \\ \hline
        $150cm \leq distance < 200cm$  & $50ms$            & $1450ms$          & $1500ms$  \\ \hline
        $100cm \leq distance < 150cm$  & $50ms$            & $ 950ms$          & $1000ms$  \\ \hline
        $ 50cm \leq distance < 100cm$  & $50ms$            & $ 700ms$          & $ 750ms$  \\ \hline
        $ 25cm \leq distance <  50cm$  & $50ms$            & $ 450ms$          & $ 500ms$  \\ \hline
        $ 10cm \leq distance <  25cm$  & $50ms$            & $ 200ms$          & $ 250ms$  \\ \hline
        $distance < 10cm$           & $50ms$            & $  75ms$          & $ 125ms$  \\ \hline\hline
    \end{tabular}
    \caption{Strobe and Chirp periods for various distances to an object}\label{tab:alarmPeriods}
\end{table}

    \section{Turn Signals} \label{sec:turnSignals}              During your lab period, the TAs will guide the class through the first modifications to the starter code that you must make, described in this section.
If you do not attend your lab period, then you must complete this section on your own.
\textbf{\textit{Except during lab time, you may }not\textit{ discuss the solutions for this section with other students.}}


\subsection{Examining the Starter Code}

\subsubsection{turn\_signals.h}

A not-uncommon pattern in embedded systems is to design the system as a state machine, or as multiple state machines.
The states are represented as the value of an enumerated type.
A couple of key advantages of writing the system as a state machine or as a collection of state machines is that it's easy to change the state based on inputs, and it's easy for the parts of the system that take action to do so based on what state the state machine is in.

In \textit{turn\_signals.h} you'll see \lstinline{enum turn_directions} which enumerates the states of a small state machine.
You'll also see the variable \lstinline{indicated_turn_direction} which will store the state of that small state machine.
In the header file, the variable is modified with the \lstinline{extern} keyword, which can be thought of as causing the \lstinline{indicated_turn_direction} variable in \textit{turn\_signals.c} to be ``public'' so that the variable can be modified by code in both \textit{turn\_signals.c} and \textit{cart\_controller.c}.


\subsubsection{turn\_signals.c}

In \textit{turn\_signals.c} you'll see:
\begin{itemize}
    \item \function{initialize_turn_signals()} where you'll put code to set variables' initial values and to configure interrupts
    \item \function{handle_button_action()} where you'll put the code that should run whenever a pushbutton is pressed
    \item \function{ISR(TIMER1_COMPA_vect)} where you'll put the code that needs to run periodically
\end{itemize}


\subsection{Detecting Button Presses}

Requirement~\ref{spec:TurnSignalControls} states that the user will indicate the direction they want to turn by pressing the appropriate pushbutton to activate the turn signal.

Read Section~5 of the Cow Pi datasheet for an overview of handling interrupts for your \developmentboard.

To detect button presses (and key presses), you will take advantage of having the NAND of the pushbuttons and the NAND of the keypad columns being input to Arduino pins D2 and D3.

Add code to \function{initialize_turn_signals()} to register \function{handle_button_action()} as your handler for button presses.
Use the \function{attachInterrupt()} and \function{digitalPinToInterrupt()} functions, as described in Section~5.3 of the datasheet.
Specify \lstinline{CHANGE} mode;
we shall explain why in Section~\ref{subsubsec:debouncing}.

Your code to respond to button presses will go in the \function{handle_button_action()} function.

\begin{lstlisting}[basicstyle=\small]
void handle_button_action(void) {
  debounce_interrupt({
    static int8_t number_of_buttons_pressed = -1;
    bool left_button_is_pressed, cowpi_right_button_is_pressed;
    while ((left_button_is_pressed = cowpi_left_button_is_pressed())
             + (right_button_is_pressed = cowpi_right_button_is_pressed())
           == number_of_buttons_pressed) {}   // busy-wait through the race condition
    /* Add code to respond to button presses */

    number_of_buttons_pressed = left_button_is_pressed + right_button_is_pressed;
  });
}
\end{lstlisting}

\subsubsection{Debouncing} \label{subsubsec:debouncing}

When you place your code that responds to button presses in \function{handle_button_action()}, place it inside the braces for the \function{debounce_interrupt()} macro.
Just as the \function{debounce_byte()} function took care of debouncing when polling mechanical inputs, the \function{debounce_interrupt()} macro defined in \textit{supplement.h} takes care of (most of) the debouncing when mechanical inputs generate interrupts.
There are two catches.

The first catch is that \function{debounce_interrupt()} only works when the interrupt handler responds to both the rising and falling edges of the input.
This is why you must register \function{handle_button_action()} using the \lstinline{CHANGE} mode.
This isn't such a terrible limitation:
Determine whether the interrupt is due to a button being pressed or released.
If pressed, then comply with Requirement~\ref{spec:TurnSignalControls};
if released, then do nothing.

The second catch is that \function{debounce_interrupt()} works by ignoring interrupts that are generated due to switchbounce.
It does not -- it cannot -- ignore switchbounce that occurs while the interrupt handler is executing.
If there were only one device attached to the input pin, then switchbounce that occurs while the interrupt handler is executing would not be a problem:
if you have a \lstinline{static bool} variable to track whether the \textit{previous} interrupt was due to the button being pressed or released, then you know that the \textit{current} interrupt is due to the button moving in the other direction.

As it is, there are two pushbuttons that could cause the interrupt, and to respond correctly, you must determine \textit{which} pushbutton was pressed by examining the inputs.
There is, however, a race condition.
Consider this scenario: the user presses a button, closing the contacts, triggering an interrupt, and the interrupt handler launches.
Before reaching the code that reads the inputs to determine which button was pressed, the switch bounces, temporarily re-opening the contacts;
when the code reads the inputs, neither button appears to be pressed.
The fix is to introduce a loop that iterates until a sensible input reading is found.
This busy-wait loop is already in the starter code.

\subsubsection{Design the Interrupt Handler}

After the busy-wait loop terminates, use the \lstinline{left_button_is_pressed} and \\ \lstinline{cowpi_right_button_is_pressed} variables to assign the appropriate state to \\ \lstinline{indicated_turn_direction}.
The code in the pushbutton interrupt handler should \textit{only} make the changes that are due to changes in the pushbuttons' positions.


\subsection{Detecting Time}

Read Section~5.2 and Section~6 of the Cow Pi datasheet to familiarize yourself with timer interrupts on the \developmentboard.
Note that Sections~6.3--6.5 are generally duplicates of each other except for their tables and a warning at the start of Section~6.3;
for you now can read Sections 6.1, 6.2, 6.4, and 6.6.
Later, you can re-visit Sections~6.3 and 6.5.

You must use a timer interrupt, along with the turn signal's state, to implement Requirement~\ref{spec:TurnSignals}.
Requirement~\ref{spec:blinkRate} identifies the timer period that you need.
Use the equation in the datasheet's Section~6.2 to iterate through possible comparison values and prescalers until you arrive at a prescaler that is available for Timer1 and a comparison value that fits in the available counter bits for that timer.
When you have done so, add code to \function{initialize_turn_signals()}:
\begin{itemize}
    \item Create a \lstinline{cowpi_timer16bit_t} pointer (in the remaining steps, we will refer to this pointer as \lstinline{timer})
    \item Use the address offset listed in the datasheet's Section~6.4 to assign the appropriate address to \lstinline{timer}
    \item Use the datasheet's Tables~9--10 to select the WGM bits for ``Clear on Timer Compare'' mode and the CS bits for your prescaler;
        use these bits and Table~8 to generate a bit vector and assign that bit vector to \lstinline{timer->control}
    \item Subtract 1 from the comparison value and assign that to \lstinline{timer->compareA}
    \item Create a \lstinline{uint8_t} pointer (which we will refer to as \lstinline{timer_interrupt_masks}), and use the address offset listed in the datasheet's Section~6.6 to assign the appropriate address to \lstinline{timer_interrupt_masks}
    \item Treating the pointer as an array, index the \lstinline{timer_interrupt_masks} array with the timer number (1), and use the datasheet's Table~14 to assign the appropriate bit vector to enable the \texttt{TIMER\textit{n}\_COMPA\_vect} interrupt, where $n$ is the timer number
%    \item Create an ISR for the \lstinline{TIMER1_COMPA_vect} interrupt using the \function{ISR()} macro as described in Section~5.2 of the datasheet
\end{itemize}

You can now add code to the ISR for \lstinline{TIMER1_COMPA_vect} to implement Requirement~\ref{spec:blinkRate}.
Based on \lstinline{indicated_turn_direction}'s state, either the left LED should blink, or the right LED should blink, or neither should blink.
\ifbool{offerdecompositionhints}{
    Introducing a \lstinline{static} variable to your ISR may be helpful.
}{}


\vspace{1cm}

You are now ready to complete the remainder of this assignment on your own.
Reminders:
\begin{itemize}
    \item Declare any global variables you use as \lstinline{volatile}.
    \item \collaborationrules
\end{itemize}

    \section{Controlling the Cart} \label{sec:cartController}   \subsection{Examining the Starter Code}

\subsubsection{Starter Code Functions and Variables}

In \textit{cart\_controller.c} you'll see:
\begin{itemize}
    \item \function{initialize_cart_system()} where you'll put code to set variables' initial values and to configure interrupts
    \item \function{control_cart()} where you'll put code that should repeated execute as part of the system's main loop
    \item \function{handle_key_action()} where you'll put the code that should run whenever a key on the keypad is pressed or released
    \item \function{ISR(TIMER2_COMPA_vect)} where you'll put the code that needs to run periodically
\end{itemize}

You may, of course, add other helper functions to \textit{cart\_controller.c} as you see fit.

You will also see a global \lstinline{volatile} variable, \lstinline{display_needs_refreshed} that you can use as part of implementing Requirement~\ref{spec:displayRefreshRate}.
There are also two \lstinline{const} variables, \lstinline{MU} and \lstinline{DEGREE} that can be used if you wish to include `\textmu' or `\textdegree' as part of your display (you are  not required to use these characters, but they are available if you wish to do so).

Recall also that \textit{turn\_signals.h} has an \lstinline{extern} variable, \lstinline{indicated_turn_direction} that is shared between \textit{turn\_signals.c} and \textit{cart\_controller.c}.

You \textit{will} need to add additional global variables to pass information between \function{control_cart()}, \function{handle_key_action()}, and \function{ISR(TIMER2_COMPA_vect)}.
Be sure to declare any global variables that you introduce as \lstinline{volatile}.

\subsubsection{Determining Your Display Module's ROM Code}

You can skip this subsection if you will not use \lstinline{MU} or \lstinline{DEGREE}.

The LCD1602's datasheet\footnote{
    \url{https://www.sparkfun.com/datasheets/LCD/HD44780.pdf}
} describes the character sets for two possible ROMs, coded ``A00'' and ``A02''.\footnote{
    No, I don't know what happened to ROM A01.
}
Both character sets include the ASCII printable characters, but they have different characters for the non-ASCII values.
Initially, the \function{initialize_cart_system()} function has this code:
\begin{lstlisting}[basicstyle=\small]
void initialize_cart_system(void) {
  display_needs_refreshed = true;
  /* USE THIS CODE TO DETERMINE WHICH ROM CODE YOUR DISPLAY MODULE HAS, AND
     SET MU and DEGREE ACCORDINGLY ON LINE 55-61. THEN COMMENT-OUT THE NEXT
     EIGHT LINES AFTER THIS COMMENT */
  char label_string[] = "mu        degree";
  char candidates_string[17] = {0};
  if (display_needs_refreshed) {
    sprintf(candidates_string, "0:%c 2:%c  0:%c 2:%c", (char)0xE4, (char)0xB5, (char)0xDF, (char)0xB0);
    display_needs_refreshed = false;
    display_string(label_string, TOP_ROW);
    display_string(candidates_string, BOTTOM_ROW);
  }
  /* END OF CODE TO DETERMINE YOUR DISPLAY MODULE'S ROM CODE */

}
\end{lstlisting}
Depending on which ROM your display module has, your display module will initially show either

ROM Code 0 \\
\texttt{
    \colorbox{LightGreen}{mu\phantom{xxxxxxxxx}degree} \vspace{-1mm}\\
    \colorbox{LightGreen}{0:\textmu\phantom{x}2:\begin{CJK*}{UTF8}{goth}ォ\end{CJK*}\phantom{xx}0:\textdegree\phantom{x}2:-}
} \\
or

ROM Code 2 \\
\texttt{
    \colorbox{LightGreen}{mu\phantom{xxxxxxxxx}degree} \vspace{-1mm}\\
    \colorbox{LightGreen}{0:{\"a}\phantom{x}2:\textmu\phantom{\small{xxx}}0:{\ss}\phantom{x}2:\textdegree}
}

If you have the display corresponding to ROM Code 0, then make sure that lines~57--58 are uncommented and that lines~60--61 are commented-out:
\begin{lstlisting}[numberstyle=\color{gray}, numbers=left, firstnumber=55]
// A couple of handy non-ASCII characters if you want them for your display
// ROM CODE 0
const char MU = (char)0xE4;
const char DEGREE = (char)0xDF;
// ROM CODE 2
// const char MU = (char)0xB5;
// const char DEGREE = (char)0xB0;
\end{lstlisting}
On the other hand, if you have the display corresponding to ROM Code 2, then make sure that lines~57--58 are commented-out and that lines~60--61 are uncommented:
\begin{lstlisting}[numberstyle=\color{gray}, numbers=left, firstnumber=55]
// A couple of handy non-ASCII characters if you want them for your display
// ROM CODE 0
// const char MU = (char)0xE4;
// const char DEGREE = (char)0xDF;
// ROM CODE 2
const char MU = (char)0xB5;
const char DEGREE = (char)0xB0;
\end{lstlisting}
If you have neither, then contact the professor for the code to create custom characters for `\textmu' and `\textdegree' (if you plan to use these characters on your display).

After you have done so, comment-out the eight lines in \function{initialize_cart_system()} that generate that display.


\subsection{Detecting Key Presses and Releases}

Add code to \function{initialize_cart_system()} to register \function{handle_key_action()} as your handler for key presses and releases.
Use the \function{attachInterrupt()} and \function{digitalPinToInterrupt()} functions, as described in Section~5.3 of the datasheet.
Specify \lstinline{CHANGE} mode.

The starter code for \function{handle_key_action()} has debouncing code similar to what you saw in \function{handle_button_action()}.
After the busy-wait loop terminates, the \lstinline{key} variable will either hold a non-NUL character, which will tell you which key press triggered this interrupt,
or \lstinline{key} will hold `\\0', in which case the \lstinline{last_key} variable will tell you which key release triggered this interrupt.

\begin{lstlisting}[basicstyle=\small]
void handle_key_action(void) {
  debounce_interrupt({
    static char last_key = 0xF0;
    char key;
    while ((key = cowpi_get_keypress()) == last_key) {}   // busy-wait through the race condition
    /* Add code to respond to key presses and releases */

    last_key = key;
  });
}
\end{lstlisting}

As described in Requirements~\ref{spec:dPad} and \ref{spec:diagnosticDisplays}, your system needs to respond to key presses for some (but not all) of the keys on the numeric keypad,
and it needs to respond to key releases for only a few of the keys.
Until you design the logic for your cart controller, you might simply print a message to the console when these particular actions happen.

As a rule, you want your interrupt-handling code to run quickly.
You certainly don't want to take the time to send a message across the slow UART connection to the Arduino IDE's Serial Monitor.
When you introduce the logic for your cart controller, be sure to remove the print statements from your interrupt handler.

\subsection{Detecting Time}

You must use a timer interrupt, to implement Requirement~\ref{spec:dataRefreshRate}.
Use that same timer interrupt\ifbool{offerdecompositionhints}{, along with a \lstinline{static} variable the ISR,}{} to determine when to set the \lstinline{display_needs_refreshed} variable.

Requirement~\ref{spec:dataRefreshRate} identifies the timer period that you need.
Use the equation in the datasheet's Section~6.2 to iterate through possible comparison values and prescalers until you arrive at a prescaler that is available for Timer2 and a comparison value that fits in the available counter bits for that timer.
When you have done so, add code to \function{initialize_cart_system()}:
\begin{itemize}
    \item Create a \lstinline{cowpi_timer8bit_t} pointer (in the remaining steps, we will refer to this pointer as \lstinline{timer})
    \item Use the address offset listed in the datasheet's Section~6.5 to assign the appropriate address to \lstinline{timer}
    \item Use the datasheet's Tables~12--13 to select the WGM bits for ``Clear on Timer Compare'' mode and the CS bits for your prescaler;
    use these bits and Table~11 to generate a bit vector and assign that bit vector to \lstinline{timer->control}
    \item Subtract 1 from the comparison value and assign that to \lstinline{timer->compareA}
    \item Create a \lstinline{uint8_t} pointer (which we will refer to as \lstinline{timer_interrupt_masks}), and use the address offset listed in the datasheet's Section~6.6 to assign the appropriate address to \lstinline{timer_interrupt_masks}
    \item Treating the pointer as an array, index the \lstinline{timer_interrupt_masks} array with the timer number (2), and use the datasheet's Table~14 to assign the appropriate bit vector to enable the \texttt{TIMER\textit{n}\_COMPA\_vect} interrupt, where $n$ is the timer number
%    \item Create an ISR for the \lstinline{TIMER1_COMPA_vect} interrupt using the \function{ISR()} macro as described in Section~5.2 of the datasheet
\end{itemize}

Until you design the logic for your cart controller, you might simply leave the ISR's body empty except for the code that updates \lstinline{display_needs_refreshed}.


\subsection{Updating the Display}

As a rule, you want your interrupt-handling code to run quickly.
You certainly don't want to take the time to update the display module in your ISR\@.
Instead, include in your ISR enough code to determine when 256 nanofortnights have passed,
\ifbool{offerdecompositionhints}{
    that is, when the timer interrupt has fired 256 times,
}{}
and set the \lstinline{display_needs_refreshed} when that has happened.

Then, in the \function{control_cart()} function, find this code:
\begin{lstlisting}[basicstyle=\small]
  static char top_line[17];
  static char bottom_line[17];
  if (display_needs_refreshed) {
    display_needs_refreshed = false;
    /* Add to create the strings to be displayed */

  }
\end{lstlisting}

Place inside that \lstinline{if} statement's body any code that should run whenever the display needs to be updated.
You will probably want to use \function{sprintf()} to generate the strings to be displayed.
See Appendix~\ref{sec:formatStrings} for some discussion.

\subsection{Controlling the Cart}

Instead of hacking code together, you should first think about its design.
You might not get the design perfect initially, but you will have much less backtracking if you give thought to it first.

\subsubsection{Data Representations}

Consider, for example, Requirement~\ref{spec:importantData}.
Finding a suitable datatype for the speed should be straight-forward enough.
What about the heading and distance?
How can you maintain a precision of \textit{exactly} 0.000,001\textdegree and of \textit{exactly} 0.001~microfurlongs?

\ifbool{offerdecompositionhints}{
    An IEEE~754 binary floating point type will not allow you to maintain precisions of exactly 0.000,001\textdegree and 0.001~microfurlongs.
    An attempt to approximate those precisions will result in the display frequently showing the incorrect distance travelled -- admittedly for only 256 nanofortnights each time this happens, but you can do better.
    (Another problem is that, by default, \function{sprintf()} will not work with floating point values on microcontrollers, including the ATmega328P.
    We can change this with a compiler argument, but the Arduino IDE doesn't make it easy to do that.)
    %\paragraph{Hint} Review your notes from the part of the floating point lecture in which we discussed how decimal currency is handled. \\
    \paragraph{Hint} Review the ``Decimal Fixed Point Value Hack'' slides from the floating point lecture, and think about how to use integer types to represent 0.001~microfurlongs and 0.000,001\textdegree.
    \paragraph{Hint} There are 1,000~nanofurlongs in a microfurlong, and there are 1,000,000~microdegrees in a degree. \\
}{}

How will you keep track of whether the motor is engaged?
How will you keep track of what direction (if any) the cart is turning?
How will you keep track of which display should be displayed (see Requirement~\ref{spec:diagnosticDisplays})?

\subsubsection{Design the Displays}

Design the diagnostic displays that show the full variables you chose.
Also design the ``normal'' display that shows the speed in furlongs per fortnight, the distance in microfurlongs, and the heading in degrees.
Be sure to allow sufficient room on the ``normal'' display to show a speed up to 3,000~furlongs per fortnight, a distance up to 9,999~microfurlongs, and a heading up to 359\textdegree.

Place the code to select the display in the \function{handle_key_action()} function, and the code to generate the displays and send them to the display module in the \function{control_cart()} function.

\subsubsection{Speed}

Now think about how you will change the speed.
How will you handle the minimum and maximum speeds?
How will you handle low and high gear?
How will you handle the motor being engaged or not engaged?

What has to happen when you press the `0' key to brake?
What else has to happen when you press the brake?
What has to happen when you release the brake?
What else has to happen when you release the brake?

\ifbool{offerdecompositionhints}{
    \paragraph{Hint} The requirements state that the brake lights shall be illuminated while as the user is pressing the brake.
    You can correctly infer that once the user presses the brake, they continue to do so until they release the brake. \\
}{}

Now that you've given thought to how to control the cart's speed, write the code for accelerating, decelerating, and braking.

\subsubsection{Distance Travelled}

Now that you can change the cart's speed, you can change the distance travelled.
Add code to \function{ISR(TIMER2_COMPA_vect)} that will update the distance travelled whenever the timer interrupt fires.

\ifbool{offerdecompositionhints}{
    \paragraph{Hint} Assuming you designed the timer interrupt to fire every nanofortnight:
    \begin{align*}
        Change\_in\_distance    & = speed \times time \\
                                & = \mathtt{speed\_variable}\ \frac{furlong}{fortnight} \times 1\ nanofortnight \\
                                & = \mathtt{speed\_variable}\ nanofurlong \\
                                & = \mathtt{speed\_variable} \times 0.001\ microfurlong
    \end{align*}
}{}

\subsubsection{Turning}

While activating a turn signal indicates the intention to turn a particular direction, actually turning requires using the ``D-pad'' to indicate the turn direction and forward motion to determine how much the direction changes.
This means that the code for making turns will be spread across a couple of functions.

Don't forget that the turn signal should stop blinking when the user ends the turn.

\subsubsection{Final Review of Requirements}

Look again over the requirements in Section~\ref{sec:spec}.
Have you forgotten anything?

Check to make sure that you've removed any long-running code, such as \function{display_string()} and \function{printf()} from your interrupt handlers / ISRs.


    \section{Turn-in and Grading}                               \filesubmission.

\policyforcodethatdoesnotcompile

\latepolicy

\subsection*{Rubric}

This assignment is worth 35 points.
\begin{description}
    \rubricitem{1}{\function{is_nan()} correctly reports whether or not its argument is a number}
    \rubricitem{1}{\function{is_zero()} correctly reports whether or not its argument is zero}
    \rubricitem{1}{\function{is_infinity()} correctly reports whether or not its argument is infinite}
    \rubricitem{1}{\function{is_negative()} correctly reports whether or not its argument is negative}
    \rubricitem{1}{\function{get_754_integer()} correctly extracts the significand's implicit integer}
    \rubricitem{1}{\function{get_754_fraction()} correctly extracts the significand's fraction bits}
    \rubricitem{1}{\function{get_754_exponent()} correctly extracts the exponent}
    \rubricitem{1}{\function{decode()} correctly converts an \lstinline{ieee754_t} value into a \lstinline{unnormal_t} structure}
    \rubricitem{1}{\function{negate()} correctly changes its argument's sign}
    \rubricitem{5}{\function{add()} can add integers \& fractions, positive \& negative values, and ``large'' \& ``small'' numbers}
    \rubricitem{1}{The identity and commutative properties hold for \function{add()}}
    \rubricitem{1}{\function{add()} provides correct answers for its special cases}
    \rubricitem{5}{\function{multiply()} can multiply integers \& fractions, positive \& negative values, and ``large'' \& ``small'' numbers}
    \rubricitem{2}{The identity, zero, and commutative properties hold for \function{multiply()}}
    \rubricitem{1}{\function{multiply()} provides correct answers for its special cases}
    \rubricitem{1}{\function{divide()} provides correct answers for its special cases}
    \rubricitem{1}{\function{divide()} can divide when the divisor is of the form $\pm 2^n, -126 \le n \le 127$}
    \rubricitem{1}{\function{divide()} can divide when the dividend's significand is a multiple of the divisor's significand}
    \rubricitem{1}{\function{add()} demonstrates that \function{encode()} rounds down when the truncated part of the significand is less than halfway between representable values}
    \rubricitem{1}{\function{add()} demonstrates that \function{encode()} rounds up when the truncated part of the significand is more than halfway between representable values}
    \rubricitem{2}{\function{add()} demonstrates that \function{encode()} rounds to the nearest-even when the truncated part of the significand is exactly halfway between representable values}
    \rubricitem{1}{Rounding can carry into the exponent}
    \rubricitem{1}{\function{add()} and/or \function{multiply()} demonstrate that \function{encode()} overflows to infinity}
    \rubricitem{1}{\function{add()}, \function{multiply()}, and/or \function{divide()} demonstrate that \function{encode()} gracefully underflows through subnormal numbers}
    \rubricitem{1}{\function{multiply()} and/or \function{divide()} demonstrate that \function{encode() underflows to zero}}
    \bonusitem{2}{\function{divide()} can divide arbitrary values}
\end{description}

\textbf{Penalties}
\begin{description}
    \softwareengineeringpenalties
    \item[no credit] for functions that use \lstinline{float} or \lstinline{double} variables or constants, use \lstinline{union} variables, use C's floating point operations, and/or a function you did not write
    \item[no credit] for arithmetic functions, if \function{decode()} and/or \function{encode()}  use \lstinline{float} or \lstinline{double} variables or constants, use \lstinline{union} variables, use C's floating point operations, and/or a function you did not write
\end{description}


    \section*{Epilogue}                                         \LauraDern

    \textit{To be continued...}

%    \newpage\appendix
    \appendix

    \section{Appendix: Lab Checkoff}                            You are not required to have your assignment checked-off by a TA or the professor.
If you do not do so, then we will perform a functional check ourselves.
In the interest of making grading go faster, we are offering a small bonus to get your assignment checked-off at the start of your scheduled lab time immediately after it is due.
Because checking off all students during lab would take up most of the lab time, we are offering a slightly larger bonus if you complete your assignment early and get it checked-off by a TA or the professor during office hours.

\begin{enumerate}
    \precheckoffitem{Position the Cow Pi a little more than 1~meter from a wall, or position an upright book (or other object) a little more than 1~meter from the Cow Pi.}
    \precheckoffitem{Place both switches in the left position.}
    \precheckoffitem{Upload your code to your \developmentboard, and leave your code open in the IDE.}
    \precheckoffitem{Confirm that the system detects the wall (or book or other object) and not something closer (such as a computer or the table surface).}

    \checkoffitem{Show and explain to the TA how your code generates a tone with a frequency of 5kHz; that is, it has a period of 200\textmu s.}
    \checkoffitem{Place the right switch in the right position, putting the system in Continuous Tone mode.
        The system generates a continuous 5kHz tone.}
    \item[] (TA, confirm that the tone is 5kHz by code inspection and by ear; confirm with the HuskerScope spectrum analyzer if you aren't sure.) \\
        \textit{+1 There is code to generate an audible tone} \\
        \textit{+1 The system continuously generates the tone when in Continuous Tone mode} \\
        \textit{+2 The audible tone has a frequency of 5kHz}

    \checkoffitem{Place the left switch in the right position, putting the system in Threshold Adjustment mode.
        The system prompts for a new threshold range. \\
        \textit{+1 The user is prompted to enter a new threshold range when the system is in Threshold Adjustment mode}}
    \checkoffitem{Enter a range of 25, using the `\#' key to indidicate that you have fininished entering the value.
        The system displays a helpful error message explaining that this is not a valid threshold range.
        The system then prompts the user for a new threshold range. \\
        \textit{+1 The user is given a helpful error message after entering an invalid threshold range} \\
        \textit{+1 The user is re-prompted to enter a threshold range after entering an invalid threshold range}}
    \checkoffitem{Enter a range of 450, using the `\#' key to indidicate that you have fininished entering the value.
        The system displays a helpful error message explaining that this is not a valid threshold range.
        The system then prompts the user for a new threshold range.}
    \checkoffitem{Enter a range of 75, using the `\#' key to indidicate that you have fininished entering the value.
        The system displays a message confirming the new threshold range. \\
        \textit{+2 Valid threshold ranges are those between 50cm and 400cm, inclusive} \\
        \textit{+1 The user is shown a confirmation message after entering a valid threshold range} \\
        \textit{+2 The user can enter a new threshold range when the system is in Threshold Adjustment mode}}

    \checkoffitem{Place the right switch in the left position, putting the system in Single Pulse mode.
        The system might indicate that no object has been detected yet; however, this is not required information before initiating a ping.}
    \checkoffitem{Show and explain to the TA how your code initiates a pulse.}
    \checkoffitem{Show and explain to the TA how your code measures the length of a pulse.}
    \checkoffitem{Show and explain to the TA how your code achieves the required precision (no greater than 1\textmu s) and accuracy (immediately detect pulse edges without waiting for code in the main loop to poll the pin).}
    \checkoffitem{Press the pushbutton to initiate a pulse.
        The right LED strobes once.
        The piezodisc does not chirp.
        The system displays the correct distance to the wall (or book or other object). \\
        \textit{+2 There is code to initiate an ultrasound pulse} \\
        \textit{+3 There is code to detect the length of the pulse} \\
        \textit{+3 The pulse's length is measured to a precision of no greater than 1\textmu s} \\
        \textit{+3 The pulse's length is measured as accurately as possible} \\
        \textit{+2 The user can request a ping when the system is in Single Pulse mode} \\
        \textit{+2 The distance to an object is correctly calculated from the pulse's length} \\
        \textit{+2 When an object is detected, the system displays the distance to the object} \\
        \textit{+2 When an object is detected in Single-Pulse mode, the system generates exactly one alarm} \\
        \textit{+2 A strobe is an illumination of the right LED for 50ms} \\
        \textit{+2 A strobe occurs for any detected object}}
    \checkoffitem{Slightly change the distance between the Cow Pi and the target object.
        Press the pushbutton to initiate a pulse.
        The LED strobes once.
        The piezodisc does not chirp.
        The system displays the new distance to the wall (or book or other object). \\
        \textit{+2 The user can request another ping when the system is in Single Pulse mode}}

    \checkoffitem{Place the left switch in the left position, putting the system in Normal Operation mode.
        The system displays the distance to the target object, and it displays an approach rate of 0cm/s.
        The LED strobes once per seccond (100ms), but the piezodisc does not chirp. \\
        \textit{+1 The switches control the mode of operation as specified} \\
        \textit{+2 When an object is detected in Normal Operation mode, the system repeatedly generates alarms}}
    \checkoffitem{Slowly move the Cow Pi closer to the wall, or slowly move the book (or other object) closer to the Cow Pi.
        As you do so, vary the rate of approach slighly to demonstrate that the rate of approach updates.
        The displayed distance changes with the decreasing distance to the target object.
        The system displays a plausible, positive rate of approach that updates at least once every second. \\
        \textit{+2 When an object is detected in Normal Operation mode the rate of approach is displayed} \\
        \textit{+2 When an object is detected in Normal Operation mode, the rate of approach is updated at least once every second}}
    \checkoffitem{As the distance between the Cow Pi and the target object decreases, note that:
        \begin{description}
            \item[When the distance falls below 100cm] the LED strobes more frequently, once every 750ms (\textthreequarters sec)
            \item[When the distance falls below 75cm] the piezodisc chirps every 750ms
            \item[When the distance falls below 50cm] the LED strobes, and the piezodisc chirps, every 500ms (\textonehalf sec)
            \item[When the distance falls below 25cm] the LED strobes, and the piezodisc chirps, every 250ms (\textonequarter sec)
            \item[When the distance falls below 10cm] the LED strobes, and the piezodisc chirps, every 125ms ($\frac{1}{8}$ sec)
        \end{description}
        \textit{+2 A chirp is an audible tone lasting 50ms} \\
        \textit{+2 When the system repeatedly generates alarms, the time between alarms is as specified}}

    \checkoffitem{Place the both switches in the right position, putting the system in Threshold Adjustment mode.
        The system prompts for a new threshold range.}
    \checkoffitem{Enter a range of 55.
        The system displays a message confirming the new threshold range.}
    \checkoffitem{Place the both switches in the left position, putting the system in Normal Operation mode.
        The system displays the distance to the target object, and it displays an approach rate of 0cm/s.}
    \checkoffitem{Slowly move the Cow Pi away from the wall, or slowly move the book (or other object) away from the Cow Pi.
        The displayed distance changes with the decreasing distance to the target object.
        The system displays a plausible, negative rate of approach.}
    \checkoffitem{As the distance between the Cow Pi and the target object decreases, the alarms become less urgent.
        Note that as the distance increases above 55cm, the piezodisc stops chirping but the LED continues to strobe. \\
        \textit{+2 A chirp occurs for a detected object that is closer than the threshold range} \\
        \textit{+2 A chirp only occurs for a detected object that is closer than the threshold range}}

    \checkoffitem{Reorient the Cow Pi, or remove the book (or other object) so that there are no in-range objects to detect.
        The system displays a message indicating that no object is detected.
        The LED does not strobe, and the piezodisc does not chirp.\\
        \textit{+3 The code correctly recognizes the that no object has been detected, if no object reflects the ultrasound pulse} \\
        \textit{+2 When there is no in-range object, the system displays a message to that effect} \\
        \textit{+2 A strobe only occurs for a detected object}}

    \checkoffitem{Show the TA any code they have not yet examined. \\
        \textit{+1 The code is clean, well-organized, has good variable and function names, and is otherwise understandable}}
\end{enumerate}

\end{document}





%Please familiarize yourself with the entire assignment before beginning.
%Section~\ref{sec:FunctionalSpecification} has the functional specification of
%the system system you will develop. Section~\ref{sec:Constraints} describes
%implementation constraints. Section~\ref{sec:Paradigm} relates interrupt-driven
%I/O to event-driven programming (which you should have learned about in \cstwo).
%Section~\ref{sec:GettingStarted} identifies code that you can reuse from PollingLab.
%Section~\ref{sec:TimerInterrupts} will guide you through configuring a hardware
%timer and responding to interrupts from that timer.
%Section~\ref{sec:ExternalInterrupts} will guide you through registering and writing
%interrupt service routines for external interrupts. Finally,
%Sections~\ref{sec:ConversionTimeout} and \ref{sec:BuildingAndConverting} offers
%suggestions for using your interrupt service routines to implement the system's
%specifications.












%\section*{Turn-in and Grading}\addcontentsline{toc}{section}{Rubric}
%
%When you have completed this assignment, upload \textit{InterruptLab.ino} to
%\filesubmission.
%
%This assignment is worth 40 points. \\
%
%Rubric:
%\begin{description}
%\item[Pushbutton Interrupts]
%\rubricitem{4}{Pushbutton presses and releases are detected with an external
%    interrupt.}
%\rubricitem{2}{The pushbuttons' interrupt handler determines which button was
%    released.}
%\rubricitem{2}{The pushbuttons' interrupt handler determines whether the left
%    button has been double-clicked.}
%\item[Matrix Keypad Interrupts]
%\rubricitem{4}{Matrix keypad presses are detected with an external interrupt.}
%\rubricitem{2}{The keypad's interrupt handler determines which key was pressed.}
%\item[Timer Interrupts]
%\rubricitem{2}{Timer interrupts for Timer1 or Timer2 are enabled and handled.}
%\rubricitem{4}{The timer interrupts configured such that, when combined with
%    other logic, the software is able to determine when exactly 7.5~seconds or
%    exactly 20~seconds have passed since the left button was double-clicked.}
%\item[Conversion Timeout]
%\rubricitem{2}{The system is able to determine when exactly 7.5~seconds have
%    passed since the left button was double-clicked.}
%\rubricitem{2}{The system is able to determine when exactly 20~seconds have
%    passed since the left button was double-clicked.}
%\item[Building Numbers]
%\rubricitem{3}{Builds a value consistent with
%    requirements~\ref{spec:initallyZero}--\ref{spec:noLeadingZeroes}.
%    (\textthreequarters\ point for each of: displaying correct positive decimal
%    values, displaying correct negative decimal values, displaying correct
%    positive hexadecimal values, displaying correct negative hexadecimal
%    values)}
%\rubricitem{1}{Negates value when left pushbutton is released. (\textonequarter\
%    point for each of: displaying correct positive decimal values, displaying
%    correct negative decimal values, displaying correct positive hexadecimal
%    values, displaying correct negative hexadecimal values)}
%\rubricitem{1}{Releasing the right pushbutton clears all digits on the Display
%    Module except the least-significant digit, which displays {\dviiseg 0}.}
%\rubricitem{1}{Detects and displays correct message when the number being built
%    is too big. (\textonehalf\ point for detecting too-big numbers;
%    \textonequarter\ point for no false detections; \textonequarter\ point for
%    displaying the correct message)}
%\item[Converting Numbers]
%\rubricitem{2}{Converts from decimal to hexadecimal when the left button is
%    double-clicked. (1 point for correctly-displayed positive numbers; 1 point
%    for correctly-displayed negative numbers)}
%\rubricitem{2}{Converts from hexadecimal to decimal when the left button is
%    double-clicked. (1 point for correctly-displayed positive numbers; 1 point
%    for correctly-displayed negative numbers)}
%\rubricitem{1}{Displays the number in its original number base 7.5~seconds or
%    20~seconds (as appropriate) after displaying the converted number.}
%\rubricitem{1}{Detects and displays correct message when the number being
%    converted to the other number base is too great for the new number base.}
%\rubricitem{1}{Displays the number in its original number base 7.5~seconds or
%    20~seconds (as appropriate) after displaying the error message.}
%\rubricitem{1}{The external LED illuminates while and only while the converted
%    number (or the convered error message) is displayed.}
%\item[Other Requirements]
%\rubricitem{1}{The number being built does not change when a button is pressed.}
%\rubricitem{1}{If the left button is double-clicked, then the number being built
%    does not change when the button is released.}
%\bonusitem{2}{Get assignment checked-off by TA or professor during office hours
%    before it is due. (You cannot get both bonuses.)}
%\bonusitem{1}{Get assignment checked-off by TA at \textit{start} of your
%    scheduled lab immediately after it is due. (Your code must be uploaded to
%    \filesubmission\ \textit{before} it is due. You cannot get both bonuses.)}
%\item[Penalties]
%\penaltyitem{0}{No penalty needed for ``Pushbutton Interrupts,'' ``Matrix Keypad
%    Interrupts,'' or ``Timer Interrupts'' portions of rubric since credit for
%    these parts of the rubric is not possible without complying with the
%    constraints in Section~\ref{sec:Constraints}.}
%\penaltyitem{4}{Code associated with conversion timeout (other than that
%    covered in the first penalty item) relies on code that violates the
%    constraints in Section~\ref{sec:Constraints}.}
%\penaltyitem{6}{Code associated with building numbers (other than that
%    covered in the first penalty item) relies on code that violates the
%    constraints in Section~\ref{sec:Constraints}.}
%\penaltyitem{8}{Code associated with converting numbers (other than that
%    covered in the first two penalty items) relies on code that violates the
%    constraints in Section~\ref{sec:Constraints}.}
%\penaltyitem{2}{Code associated with certain button actions having no effect
%    (other than that covered in the first penalty item) relies on code that
%    violates the constraints in Section~\ref{sec:Constraints}.}
%\spaghetticodepenalties{1}
%\end{description}
%%
%
%
%
%
%\section*{Appendix: Lab Checkoff}\addcontentsline{toc}{section}{Lab Checkoff}
%
%You are not required to have your assignment checked-off by a TA or the
%professor. If you do not do so, then we will perform a functional check
%ourselves. In the interest of making grading go faster, we are offering a small
%bonus to get your assignment checked-off at the start of your scheduled lab
%time immediately after it is due. Because checking off all students during lab
%would take up most of the lab time, we are offering a slightly larger bonus if
%you complete your assignment early and get it checked-off by a TA or the
%professor during office hours.
%
%\begin{description}
%\precheckoffitem{Establish that the code you are demonstrating is the code
%    you submitted to to \filesubmission.}
%    \begin{itemize}
%    \item If you are getting checked-off during lab time, show the TA that the
%        file was submitted before it was due.
%    \item Download the file into your PollingLab directory. If necessary,
%        rename it to \textit{InterruptLab.ino}.
%    \end{itemize}
%\precheckoffitem{Upload \textit{InterruptLab.ino} to your \nano\ and open the
%    Serial Monitor.}
%\end{description}
%
%\begin{enumerate}
%\checkoffitem{Show in your code that your \function{loop()} function has no
%    executable code, and that there is not an infinite loop in the
%    \function{setup()} fuction, nor in any functions called by
%    \function{setup()}.} \\
%    \textit{establishing this now allows us to conclude that the observed
%    functionality is in the ISRs or their helper functions}
%\checkoffitem{Show in your code that you registered interrupt handlers for
%    the pushbuttons and for the matrix keypad.} \\
%    \textit{make note of the interrupt condition}
%\checkoffitem{Show in your code your interrupt handlers for the pushbuttons and
%    for the matrix keypad.} \\
%    \textit{+4 pushbutton presses and releases are detected with an external
%    interrupt \\
%    \phantom{xxx}(if a single handler is registered for pushbutton CHANGEs, or \\
%    \phantom{xxx}if handlers are registered for both FALLING and RISING
%    pushbuttons)} \\
%    \textit{+4 matrix keypad presses are detected with an external interrupt}
%\checkoffitem{Show and explain how you determined the comparison value and
%    prescaler for your timer interrupts.}
%\checkoffitem{Show in your code that you had configured the timer to use
%    those comparison and prescaler values and that you had configured the
%    correct timer mode.} \\
%    \textit{+4 timer interrupts are configured correctly}
%\checkoffitem{Show in your code that you enabled the correct timer.}
%\checkoffitem{Show in your code that you created an ISR for the correct
%    timer.} \\
%    \textit{+2 timer interrupts are enabled and handled}
%\checkoffitem{Place the right switch in the left (decimal) position.}
%\checkoffitem{Press 1. One of two things will happen:}
%    \begin{itemize}
%    \item 1 appears on the Serial Monitor.
%    \item {\dviiseg 1} appears on the display module.
%    \end{itemize}
%    \textit{+2 keypad ISR determines which key was pressed}
%\checkoffitem{Press the left pushbutton. The Serial Monitor may indicate that
%    the left pushbutton is DOWN and the right pushbutton is UP, but there is
%    no other observable behavior.}
%\checkoffitem{Release the left pushbutton. One or both of two things will
%    happen:}
%    \begin{itemize}
%    \item The Serial Monitor indicates that both pushbuttons are UP.
%    \item {\dviiseg -1} appears on the display module.
%    \end{itemize}
%\checkoffitem{Press the right pushbutton. The Serial Monitor may indicate that
%    the right pushbutton is DOWN and the left pushbutton is UP, but there is
%    no other observable behavior.} \\
%    \textit{+1 The number being build does not change when a button is pressed.}
%\checkoffitem{Release right pushbutton. One or both of two things will
%    happen:}
%    \begin{itemize}
%    \item The Serial Monitor indicates that both pushbuttons are UP.
%    \item {\dviiseg 0} appears on the display module.
%    \end{itemize}
%    \textit{+2 pushbutton ISR determines which button was released}
%\checkoffitem{Place the right switch in the right (hexadecimal) position. Place
%    the right switch in the left (7.5~seconds) position.}
%\checkoffitem{Press A. {\dviiseg A} may appear on the display module.}
%\checkoffitem{Double-click the left pushbutton. One or more of three things will
%    happen:}
%    \begin{itemize}
%    \item The Serial Monitor indicates that a double-click occurred.
%    \item The display will update to {\dviiseg 10} or an attempt to do so, such
%        as {\dviiseg 0} or {\dviiseg 1}.
%    \item The external LED illuminates.
%    \end{itemize}
%    \textit{+2 pushbutton ISR determines whether the left button is
%    double-clicked} \\
%\checkoffitem{Wait 7.5~seconds. One or more of three things will happen:}
%    \begin{itemize}
%    \item The Serial Monitor indicates that 7.5~seconds have elapsed.
%    \item The display will revert to {\dviiseg A} or an attempt to do so.
%    \item The external LED deluminates.
%    \end{itemize}
%    \textit{+2 system determines when 7.5~seconds have passed} \\
%    \textit{+1 external LED illuminates while and only while converted number is
%    displayed} \\
%% \checkoffitem{Release both buttons. The Serial Monitor may briefly indicate that
%%     one button is DOWN and the other is UP, and it may indicate that both
%%     buttons are UP, but there is no other observable behavior.} \\
%    \textit{+1 the number being built does not change when the left button is
%    released during a double-click}
%\checkoffitem{Place the left switch in the left (20~seconds) position.}
%\checkoffitem{Double-click the left pushbutton. One or more of the
%    previously-mentioned indications of simultaneous button presses happen.}
%\checkoffitem{Wait 20~seconds. One or more of the previously-mentioned
%    indications of conversion timeout happen.} \\
%    \textit{+2 system determines when 20~seconds have passed}
%\checkoffitem{Place the left switch in the right (7.5~seconds) position.}
%\checkoffitem{Press and release the right pushbutton. The output is {\dviiseg 0}
%    on the display module.}
%\checkoffitem{Enter the value 0x123A. The output is {\dviiseg\phantom{8888}123A}
%    on the display module.} \\
%    \textit{+\textthreequarters\ displays correct positive hexadecimal values}
%\checkoffitem{Press the and release the left pushbutton. The output is
%    {\dviiseg FFFFEDC6} on the display module.} \\
%    \textit{+\textonequarter\ correctly negates positive hexadecimal values}
%\checkoffitem{Press and release the right pushbutton. The output is {\dviiseg 0}
%    on the display module.}
%    \textit{+1 releasing the right pushbutton clears all digits and then
%        displays {\dviiseg 0}}
%\checkoffitem{Enter the value 0xB654789C. The output is {\dviiseg B654789C}
%    on the display module.} \\
%    \textit{+\textthreequarters\ displays correct negative hexadecimal values}
%\checkoffitem{Press the and release the left pushbutton. The output is
%    {\dviiseg 49AB8764} on the display module.} \\
%    \textit{+\textonequarter\ correctly negates negative hexadecimal values}
%\checkoffitem{Place the right switch in the left (decimal) position.}
%\checkoffitem{Press and release the right pushbutton. The output is {\dviiseg 0}
%    on the display module.}
%\checkoffitem{Enter the value 1478. The output is {\dviiseg\phantom{8888}1478}
%    on the display module.} \\
%    \textit{+\textthreequarters\ displays correct positive decimal values}
%\checkoffitem{Press the and release the left pushbutton. The output is
%    {\dviiseg\phantom{888}-1478} on the display module.} \\
%    \textit{+\textonequarter\ correctly negates positive decimal values}
%\checkoffitem{Press and release the right pushbutton. The output is {\dviiseg 0}
%    on the display module.}
%\checkoffitem{Enter the value -5236. The output is {\dviiseg\phantom{888}-5236}
%    on the display module.} \\
%    \textit{+\textthreequarters\ displays correct negative decimal values}
%\checkoffitem{Press the and release the left pushbutton. The output is
%    {\dviiseg\phantom{8888}5236} on the display module.} \\
%    \textit{+\textonequarter\ correctly negates negative decimal values}
%\checkoffitem{Press and release the right pushbutton. The output is {\dviiseg 0}
%    on the display module.}
%\checkoffitem{Enter the value 12369874. The output is {\dviiseg 12369874} on
%    the display module.} \\
%    \textit{+\textonequarter\ no false ``too big'' detections}
%\checkoffitem{Press 0. The output is {\dviiseg\phantom{8}too big} on the display
%    module.} \\
%    \textit{+\textonehalf\ detects that a number is ``too big''} \\
%    \textit{+\textonequarter\ displays correct error message}
%\checkoffitem{Press and release the right pushbutton. The output is {\dviiseg 0}
%    on the display module.}
%\checkoffitem{Enter the value 1234. The output is {\dviiseg\phantom{8888}1234}
%    on the display module.}
%\checkoffitem{Double-click the left pushbutton. The output is
%    {\dviiseg\phantom{88888}4D2} on the display module.} \\
%    \textit{+1 converts positive decimal to hexadecimal}
%\checkoffitem{Wait 7.5~seconds for the conversion to timeout. The output is
%    {\dviiseg\phantom{8888}1234} on the display module.} \\
%    \textit{+\textonequarter\ the original positive decimal number displays after
%    timeout}
%\checkoffitem{Press and release the left pushbutton. The output is
%    {\dviiseg\phantom{888}-1234} on the display module.}
%\checkoffitem{Double-click the left pushbutton. The output is
%    {\dviiseg FFFFFB2E} on the display module.} \\
%    \textit{+1 converts negative decimal to hexadecimal}
%\checkoffitem{Wait 7.5~seconds for the conversion to timeout. The output is
%    {\dviiseg\phantom{888}-1234} on the display module.} \\
%    \textit{+\textonequarter\ the original negative decimal number displays after
%    timeout}
%\checkoffitem{Press and release the right pushbutton. The output is {\dviiseg 0}
%    on the display module.}
%\checkoffitem{Place the right switch in the right (hexadecimal) position.}
%\checkoffitem{Enter the value 0x56B7. The output is
%    {\dviiseg\phantom{8888}56B7} on the display module.}
%\checkoffitem{Double-click the left pushbutton. The output is
%    {\dviiseg\phantom{888}22199} on the display module.} \\
%    \textit{+1 converts positive hexadecimal to decimal}
%\checkoffitem{Wait 7.5~seconds for the conversion to timeout. The output is
%    {\dviiseg\phantom{8888}56B7} on the display module.} \\
%    \textit{+\textonequarter\ the original positive hexadecimal number displays
%    after timeout}
%\checkoffitem{Press and release the left pushbutton. The output is
%    {\dviiseg FFFFA949} on the display module.}
%\checkoffitem{Double-click the left pushbutton. The output is
%    {\dviiseg\phantom{88}-22199} on the display module.} \\
%    \textit{+1 converts negative hexadecimal to decimal}
%\checkoffitem{Wait 7.5~seconds for the conversion to timeout. The output is
%    {\dviiseg FFFFA949} on the display module.} \\
%    \textit{+\textonequarter\ the original negative hexadecimal number displays
%    after timeout}
%\checkoffitem{Press and release the right pushbutton. The output is {\dviiseg 0}
%    on the display module.}
%\checkoffitem{Enter the value 0x6A0,0000. The output is
%    {\dviiseg\phantom{8}6A00000} on the display module.}
%\checkoffitem{Double-click the left pushbutton. The output is
%    {\dviiseg\phantom{888}ERROR} on the display module.} \\
%    \textit{+\textonehalf\ detects that converted number is too great to
%    display} \\
%    \textit{+\textonequarter\ displays correct error message}
%\checkoffitem{Wait 7.5~seconds for the conversion to timeout. The output is
%    {\dviiseg\phantom{8}6A00000} on the display module.} \\
%    \textit{+1 the original number displays after timeout}
%\checkoffitem{Press and release the right pushbutton. The output is {\dviiseg 0}
%    on the display module.}
%\checkoffitem{Enter the value 0x5F5,E0FF. The output is
%    {\dviiseg\phantom{8}5F5E0FF} on the display module.}
%\checkoffitem{Double-click the left pushbutton. The output is
%    {\dviiseg 99999999} on the display module.} \\
%    \textit{+\textonequarter\ no false ``error'' detections}
%\end{enumerate}
%
%This concludes the demonstration of your system's functionality. The TAs will
%later examine your code for violations of the assignment's constraints. If your
%code looks like it is tailored for this checklist, the TAs may re-grade using a
%different checklist.
